\documentclass[11pt]{article}

\usepackage[margin=1in]{geometry}
\usepackage{amsmath,amssymb}
\usepackage{booktabs}
\usepackage{array}
\usepackage[strings]{underscore}
\usepackage{hyperref}
\usepackage{xcolor}
\usepackage{enumitem}
\usepackage{caption}
\usepackage{longtable}

\hypersetup{colorlinks=true, linkcolor=blue!50!black, urlcolor=blue!50!black, citecolor=blue!50!black}

\newcommand{\rmae}{r^{\text{MAE}}}
\newcommand{\rrmse}{r^{\text{RMSE}}}
\DeclareMathOperator*{\argmin}{arg\,min}

\title{\textbf{Empirical Credibility Weighting in Paid Chain-Ladder Reserving:}\\
\large Methodology, Walk-Forward Backtesting Design, and Cross-Sectional\\
Analysis of Company Size and Loss-Experience Volatility}
\author{reserveLab \\ \small Private Passenger Automobile, NAIC Schedule P (CAS Loss Reserving Database)}
\date{\today}

\begin{document}

\maketitle

\begin{abstract}
\noindent
This report documents the full methodology, backtesting design, and statistical analysis of an
empirical credibility-weighting study built on the NAIC Schedule P private-passenger-automobile
extract of the CAS Loss Reserving Database (143 companies, accident years 1998--2007, development
lags 1--10). For each company and each development lag, a blend ratio $r \in [0,1]$ between the
company's own age-to-lag development factor and the pooled industry factor is fitted by exact,
closed-form minimization of out-of-sample walk-forward backtest error, separately under mean
absolute error (MAE) and root mean squared error (RMSE) loss. The resulting fitted ratios are
then examined cross-sectionally against (i) each company's share of total industry paid losses
and (ii) a measure of each company's loss-experience volatility, to test the credibility-theoretic
hypothesis that companies with larger, more stable books of business warrant greater reliance on
their own historical development pattern. Pearson correlation between company size and fitted
credibility is statistically indistinguishable from zero ($r=0.024$, $p=0.45$), while
rank-based measures (Spearman $\rho = 0.129$, $p<0.001$; Kendall $\tau_b = 0.093$, $p<0.001$)
detect a weak but real monotonic relationship, a discrepancy attributable to heavy corner-clustering
of the fitted ratio at the boundary values $0$ and $1$ (57.2\% of reliable observations under MAE
loss). A directly measured volatility statistic --- the sample standard deviation of year-over-year
percentage change in ultimate paid losses --- is shown to correlate strongly with prediction error
(Spearman $\rho=0.479$, $p<10^{-6}$) and strongly \emph{negatively} with company size
(Spearman $\rho=-0.683$, $p<10^{-15}$), suggesting that volatility, not size per se, is the more
fundamental driver, consistent with the process-variance basis of classical credibility theory.
Volatility does not, however, predict the \emph{direction} of the fitted blend ratio with comparable
strength, a result explained by within-company variation of $r$ across development lags. Limitations
of the underlying data --- extreme size concentration, short per-lag backtest samples, and
unregularized point-estimate optimization --- are discussed, together with suggested extensions.
\end{abstract}

\tableofcontents

%==============================================================================
\section{Introduction and Objective}
%==============================================================================

The central actuarial question addressed here is a practical reserving problem: when estimating
a private-passenger-automobile insurer's ultimate losses using the paid chain-ladder method, how
much weight should be placed on that company's own historical loss-development pattern, versus
the pattern implied by the broader industry? This is a credibility-weighting question in the
classical sense (B\"uhlmann, 1967; see Appendix~\ref{app:buhlmann} for the textbook formula used
for contrast), but rather than adopting a parametric credibility model \emph{a priori}, this
project fits the blend ratio empirically, directly from out-of-sample backtest error, and then
asks whether the resulting fitted ratios behave the way credibility theory predicts they should:
larger, more stable companies should receive more credibility for their own data.

Section~\ref{sec:data} describes the data source and population. Section~\ref{sec:cl} defines the
paid chain-ladder mechanics used throughout. Section~\ref{sec:blend} defines the credibility blend.
Section~\ref{sec:walkforward} describes the walk-forward backtesting design that generates enough
independent observations per development lag to fit anything at all. Section~\ref{sec:optimization}
derives the exact closed-form/candidate-search optimizers used for the blend ratio under MAE and
RMSE loss, with a fully worked numerical example. Section~\ref{sec:casestudies} presents
illustrative single-company case studies. Section~\ref{sec:crosssection} presents the cross-sectional
correlation and decile-binned analysis of credibility versus company size. Section~\ref{sec:volatility}
introduces the volatility measure and its correlations with error, size, and the fitted ratio itself.
Section~\ref{sec:discussion} discusses how properties of the data source shape every result above.
Section~\ref{sec:limitations} lists limitations and proposed extensions.

%==============================================================================
\section{Data}
\label{sec:data}
%==============================================================================

\subsection{Source and Provenance}

The raw data (\texttt{data/raw/ppauto\_pos.csv}) is an extract of the private-passenger-automobile
line from the CAS Loss Reserving Database, itself compiled from NAIC Schedule P, Part 1 filings ---
the loss-development exhibit every U.S. property-casualty insurer or group is required to file
annually with state insurance regulators. The extract contains 13{,}250 rows and 13 raw fields per
row: \texttt{GRCODE} (NAIC group/company code), \texttt{GRNAME} (company name), \texttt{AccidentYear},
\texttt{DevelopmentYear}, \texttt{DevelopmentLag}, \texttt{IncurredLosses}, \texttt{CumPaidLoss},
\texttt{BulkLoss} (bulk~+~IBNR reserve), \texttt{EarnedPremDIR}, \texttt{EarnedPremCeded},
\texttt{EarnedPremNet}, \texttt{Single} (1 if a single legal entity, 0 if an NAIC-defined group
aggregate of pooled subsidiaries), and \texttt{PostedReserves2007}.

\subsection{Temporal Structure: a Genuine, Disclosed ``Square,'' not an Extrapolation}
\label{sec:square}

Accident years run 1998--2007 (10 accident years) and development lag runs 1--10, with
$\texttt{DevelopmentYear} = \texttt{AccidentYear} + \texttt{DevelopmentLag} - 1$. Critically, the
raw file's \texttt{DevelopmentYear} field ranges from 1998 through \textbf{2016}, not 2007: for
accident year 2007, development lag 10 corresponds to development year 2016. This means the
dataset was compiled retrospectively, using Schedule~P filings as late as calendar year 2016 to
disclose the \emph{true, realized} cumulative paid loss at full maturity (lag 10) for every
accident year 1998--2007 -- including accident year 2007, whose lag-10 outcome would not actually
have been knowable as of the December 31, 2007 valuation date used throughout this project. This
project refers to this fully-populated accident-year $\times$ development-lag matrix as the
\emph{square}, in contrast to the \emph{triangle}: the triangle is the square restricted to the
cells that would genuinely have been observable as of a given valuation year $v$ (i.e., cells with
$\text{AccidentYear} + \text{DevelopmentLag} - 1 \le v$). Every backtest and every ``actual paid''
figure used for error grading in this project uses the square's true, disclosed future value ---
never a model-based extrapolation -- which is what makes rigorous out-of-sample grading possible
at all. (Earlier project notes used the non-standard term ``rectangle''; it has been renamed
throughout the codebase to ``square'' for clarity to other actuaries.)

\subsection{Population Characteristics}
\label{sec:population}

The extract contains $143$ distinct \texttt{GRCODE} entities: $107$ are single legal entities
(\texttt{Single}=1) and $36$ are NAIC group aggregates (\texttt{Single}=0). Total direct earned
premium (summed 1998--2007) is extremely right-skewed:

\begin{table}[h]
\centering
\caption{Direct earned premium distribution across the 143 companies (1998--2007 total)}
\begin{tabular}{lr}
\toprule
Statistic & Value \\
\midrule
Median company & \$876{,}870 \\
Mean company & \$16{,}128{,}410 \\
Largest company (State Farm Mut Grp) & \$1{,}615{,}154{,}280 \\
Top 10 companies' share of total industry premium & 92.08\% \\
Top 20 companies' share of total industry premium & 95.22\% \\
\bottomrule
\end{tabular}
\end{table}

\begin{table}[h]
\centering
\caption{Ten largest companies by total direct earned premium, 1998--2007}
\begin{tabular}{lr}
\toprule
Company & Total direct earned premium \\
\midrule
State Farm Mut Grp & \$1{,}615{,}154{,}280 \\
United Services Automobile Assn Grp & \$274{,}232{,}260 \\
FL Farm Bureau Grp & \$65{,}413{,}820 \\
New Jersey Manufacturers Grp & \$45{,}496{,}140 \\
Kentucky Farm Bureau Mut Ins Grp & \$23{,}365{,}150 \\
NC Farm Bureau Ins Grp & \$21{,}186{,}780 \\
Tenn Farmers Mut & \$21{,}172{,}580 \\
Federal Ins Co Grp & \$21{,}114{,}610 \\
Old American Cty Mut Fire Ins Co & \$19{,}916{,}490 \\
Home State Cnty Mut Ins Co & \$16{,}556{,}360 \\
\bottomrule
\end{tabular}
\end{table}

At the opposite extreme, ten companies report \emph{zero} total direct earned premium over the
entire ten-year window (e.g., Lumber Ins Cos, San Antonio Reins Co, Adriatic Ins Co, Allegheny Cas
Co) while still carrying nonzero paid-loss entries in the triangle -- consistent with runoff or
reinsurance-only entities rather than active writers of new private-passenger-auto business. This
concentration is directly relevant to Sections~\ref{sec:crosssection}--\ref{sec:discussion}.

\subsection{Coverage Completeness}

Not every company has data for all ten accident years; a company that exited the line, was
acquired, or stopped filing has a shorter observed history. Of the 143 companies, $121$ (84.6\%)
have the full 10 accident years; the remaining $22$ are distributed as follows:

\begin{table}[h]
\centering
\caption{Number of companies by count of accident years with any recorded data}
\begin{tabular}{lrrrrrrrrr}
\toprule
Accident years present & 2 & 3 & 4 & 5 & 6 & 7 & 8 & 9 & 10 \\
\midrule
Number of companies & 2 & 3 & 6 & 1 & 2 & 5 & 1 & 2 & 121 \\
\bottomrule
\end{tabular}
\end{table}

This project flags companies with fewer than $\texttt{THIN\_HISTORY\_ACCIDENT\_YEARS} = 5$
accident years as \texttt{thin\_history\_flag = True} (a boolean column persisted onto every
saved results row, surfaced as an explicit on-chart warning), since both the development-factor
fit and the backtested credibility ratio for such companies rest on almost no data. Company
$41700$ (Southern Group Ind Inc), with data for only accident years 1998--1999, is the extreme
case discussed in Section~\ref{sec:thin}.

\subsection{Pipeline Coverage}

Of the 143 companies, $112$ were processed successfully end-to-end by the full estimation
pipeline. The remaining $31$ fail deterministically at the triangle-construction stage, tripping
an explicit data-validity guard (a zero denominator in an age-to-age factor -- i.e., a company
whose cumulative paid loss is exactly zero at some development lag while nonzero at an adjacent
lag, making the ratio undefined). This is treated as a legitimate data-quality guard rather than a
defect; all cross-sectional results in Sections~\ref{sec:crosssection}--\ref{sec:volatility} are
therefore drawn from the $112$ successfully processed companies (further reduced by per-analysis
reliability filters described in those sections).

%==============================================================================
\section{Paid Chain-Ladder Methodology}
\label{sec:cl}
%==============================================================================

\subsection{Notation}

Throughout, let:
\begin{itemize}[leftmargin=2.2cm]
\item[$c$] index a company (or, when noted, the pooled industry aggregate $I$);
\item[$a$] denote accident year;
\item[$v$] denote valuation year (the ``as-of'' date at which a snapshot of the data is taken);
\item[$L$] denote development lag, $L = v - a + 1$, $L \in \{1, \dots, 10\}$;
\item[$P^{(c)}_{a,L}$] denote cumulative paid loss for company $c$, accident year $a$, at
  development lag $L$ -- one cell of the triangle or square;
\item[$f^{(c)}_{L}(v)$] denote the \emph{selected} (volume-weighted) age-to-age factor from lag
  $L$ to $L+1$, computed for company $c$ using only data observable as of valuation year $v$;
\item[$g^{(c)}_{L}(v)$] denote the \emph{age-to-lag} (age-to-ultimate, with ``ultimate'' defined
  as lag 10) factor at lag $L$, computed as of valuation year $v$.
\end{itemize}
The explicit dependence on $v$ is essential: because of the walk-forward design of
Section~\ref{sec:walkforward}, the \emph{same} company and the \emph{same} development lag are
observed with \emph{different} factor values at different points in the walk, since each point
recomputes $f$ and $g$ using only the data legitimately available as of that particular $v$ (no
look-ahead). Superscript $(c)$ denotes a company-specific quantity computed from that company's own
triangle; superscript $(I)$ denotes the identical calculation performed on the \emph{industry
triangle}, defined as the cellwise sum of $P^{(c)}_{a,L}$ across all 143 companies.

\subsection{Individual and Selected (Volume-Weighted) Age-to-Age Factors}

For a triangle (company-specific or industry) observed as of valuation year $v$, the
\emph{individual} age-to-age factor for accident year $a$ at lag $L \to L+1$ is
$P_{a,L+1} / P_{a,L}$, defined whenever both cells are observed ($a + L - 1 \le v$ and
$a + L \le v$). The \emph{selected} factor used throughout this project is the
\emph{volume-weighted average} (the industry-standard ``all-year volume-weighted'' selection,
consistent with $\texttt{factor\_average} = \texttt{"volume"}$ in the codebase), not a simple
average of individual factors:
\begin{equation}
f_{L}(v) \;=\; \frac{\displaystyle\sum_{a \,:\, a+L \le v} P_{a,L+1}}{\displaystyle\sum_{a \,:\, a+L \le v} P_{a,L}}
\label{eq:selected-factor}
\end{equation}
i.e., the sum of all next-lag cumulative paid amounts across every accident year that has reached
lag $L+1$ as of $v$, divided by the sum of the corresponding current-lag amounts. This weights
larger accident years more heavily than smaller ones, which is the standard justification for
volume weighting over a simple average of ratios.

\subsection{Age-to-Lag (Age-to-Ultimate) Factors}

The age-to-lag factor at lag $L$ is the cumulative product of selected factors from $L$ up to the
final lag (10, treated as ``ultimate'' within this 10-lag window):
\begin{equation}
g_{L}(v) \;=\; \prod_{j=L}^{9} f_{j}(v), \qquad g_{10}(v) = 1.
\label{eq:age-to-lag}
\end{equation}
Equivalently, and as implemented, $g_L(v)$ is built by a backward recursion starting from
$g_{10}(v)=1$ and repeatedly multiplying in $f_{L}(v)$ while stepping $L$ down from 9 to 1.

\subsection{Reserve and Ultimate Estimation}

For accident year $a$ observed as of valuation year $v$, let $L = v-a+1$ be its currently observed
lag and $P_{a,L}$ its last observed cumulative paid loss. The chain-ladder point estimate of
ultimate (lag-10) cumulative paid loss, and the implied reserve, are:
\begin{align}
\widehat{P}_{a,10}(v) &= P_{a,L} \cdot g_{L}(v) \label{eq:estimate}\\
\widehat{\text{Reserve}}_a(v) &= \widehat{P}_{a,10}(v) - P_{a,L}. \label{eq:reserve}
\end{align}

\subsection{Two Pattern Sources: Company versus Industry}

Every quantity above is computed twice per accident year: once using $g_L^{(c)}(v)$, the
development factors from company $c$'s own triangle (\texttt{pattern\_source = "company"}), and
once using $g_L^{(I)}(v)$, the identical calculation applied to the pooled industry triangle
(\texttt{pattern\_source = "industry"}). Both use the \emph{same} last observed paid amount
$P_{a,L}$ (the company's own), so the two estimates differ only through the choice of development
factor.

%==============================================================================
\section{Credibility-Weighted Blending}
\label{sec:blend}
%==============================================================================

For each company $c$ and development lag $L$, a blend ratio $r_{c,L} \in [0,1]$ is fit (Section
\ref{sec:optimization}) and used to form a credibility-weighted age-to-lag factor:
\begin{equation}
g_{L}^{(c,\text{blend})} \;=\; r_{c,L}\, g_{L}^{(c)} \;+\; \bigl(1-r_{c,L}\bigr)\, g_{L}^{(I)}.
\label{eq:blend-factor}
\end{equation}
This is the standard linear credibility form, with $r_{c,L}$ playing the role of the classical
credibility factor $Z$ (Appendix~\ref{app:buhlmann}): $r_{c,L}=1$ means full credibility is given
to the company's own pattern, $r_{c,L}=0$ means the company's own pattern is disregarded entirely
in favor of the industry pattern, and intermediate values are a linear interpolation. The blended
ultimate estimate and reserve follow directly by substituting Eq.~\eqref{eq:blend-factor} into
Eqs.~\eqref{eq:estimate}--\eqref{eq:reserve}:
\begin{equation}
\widehat{P}_{a,10}^{\text{blend}} = P_{a,L} \cdot g_{L}^{(c,\text{blend})}, \qquad
\widehat{\text{Reserve}}_a^{\text{blend}} = \widehat{P}_{a,10}^{\text{blend}} - P_{a,L}.
\label{eq:blend-estimate}
\end{equation}
Because $r_{c,L}$ is fit once per development lag from historical (out-of-sample) backtest data
(Section~\ref{sec:walkforward}) and then \emph{applied}, unchanged, to the target valuation year
(2007), Eq.~\eqref{eq:blend-estimate} at the target year is a genuine out-of-sample forecast, not a
refit -- a point of some consequence, returned to in Section~\ref{sec:casestudies}.
Two ratios are fit and carried in parallel throughout -- $\rmae_{c,L}$ (MAE-optimal) and
$\rrmse_{c,L}$ (RMSE-optimal) -- since the two loss functions can, and empirically do, disagree.

%==============================================================================
\section{Walk-Forward (Rolling-Origin) Backtesting Design}
\label{sec:walkforward}
%==============================================================================

\subsection{Motivation: Degeneracy of a Single Valuation Date}

At a single, fixed valuation year $v_0$ (e.g., 2007), the mapping from accident year to observed
lag, $L = v_0 - a + 1$, is a bijection: each accident year lands at exactly one lag, and each lag
is populated by exactly one accident year. Fitting a blend ratio $r_L$ per lag from a single
valuation date therefore means fitting one free parameter from exactly one data point per lag --
an exact, degenerate fit with zero residual degrees of freedom and no ability to distinguish
signal from noise.

\subsection{Construction: Walking the Valuation Date Forward}

To obtain multiple independent observations per lag, the valuation date is walked forward one
year at a time, from $v = a_0 + 1$ (where $a_0$ is the company's first observed accident year --
the earliest valuation date at which two development lags exist to form even one age-to-age
factor) through the target valuation year $v^\star = 2007$, inclusive. At each step $v$ in this
walk:
\begin{enumerate}
\item The triangle (company and industry) is rebuilt using \emph{only} cells with
  $a + L - 1 \le v$ -- i.e., properly censored as of that historical valuation date, with no
  information from calendar years after $v$ used in the factor calculation;
\item Selected factors $f_{L}(v)$ and age-to-lag factors $g_L(v)$ are recomputed from this
  $v$-censored triangle (Eqs.~\ref{eq:selected-factor}--\ref{eq:age-to-lag});
\item Every accident year $a$ observable as of $v$ contributes one row, tagged with
  $L = v-a+1$ and with $v$ itself, to a pooled table;
\item The row's \texttt{actual\_paid} field is populated from the \emph{square} (the true,
  eventually-disclosed lag-10 value for that accident year -- Section~\ref{sec:square}), used
  purely for later grading and never as a model input.
\end{enumerate}
Only the terminal step ($v = v^\star$) is persisted to disk as a deliverable; every earlier step
is transient backtest scratch work. This produces, for the \emph{same} accident year, one
observation per valuation year it lived through (at a different lag each time), and for the
\emph{same} lag, one observation per valuation year in the walk (from a different accident year
each time) -- exactly the cross-classification structure needed to fit a per-lag ratio from
multiple genuinely independent historical points.

\subsection{Concrete Worked Mapping}

For company 43, development lag $L=6$ is fed by accident years $\{1998,1999,2000,2001,2002\}$,
each paired with the valuation year at which it first reached lag 6 ($v = a+5$):
\begin{equation}
(a,v) \in \bigl\{(1998,2003),\,(1999,2004),\,(2000,2005),\,(2001,2006),\,(2002,2007)\bigr\}
\end{equation}
-- five independent $(a,v)$ pairs, five independently-computed pairs of company/industry factors
(each computed from that pair's own $v$-censored triangle), and five independent actual-vs-estimate
residuals feeding the fit for $r_{43,6}$.

\subsection{Achieved Sample Sizes}

Because the walk has a fixed length (bounded by the number of valuation years between a company's
first usable valuation date and 2007) and a fixed number of lags, observation counts shrink near
the boundary lags (very early or very late lags have fewer valid $(a,v)$ pairs within the walk).
Pooling every $(\text{company}, L)$ combination across all 112 successfully processed companies
gives $1{,}120$ rows in total, with the following distribution of \texttt{num\_observations} (the
count of $(a,v)$ pairs backing that row's fit):

\begin{table}[h]
\centering
\caption{Distribution of backtest observation counts across all 1{,}120 pooled (company, lag) rows}
\begin{tabular}{lrrrrrrrrr}
\toprule
num\_observations & 1 & 2 & 3 & 4 & 5 & 6 & 7 & 8 & 9 \\
\midrule
Count of rows & 114 & 127 & 120 & 130 & 109 & 108 & 112 & 102 & 198 \\
\bottomrule
\end{tabular}
\end{table}

Rows with \texttt{num\_observations} $=1$ are, by construction, an exact fit to a single historical
point and are treated as unreliable; all cross-sectional analyses in
Sections~\ref{sec:crosssection}--\ref{sec:volatility} exclude them, retaining $1{,}120-114=1{,}006$
rows.

%==============================================================================
\section{Exact Optimization of the Blend Ratio}
\label{sec:optimization}
%==============================================================================

\subsection{Error as an Affine Function of $r$}

Fix a company $c$ and development lag $L$, with pooled backtest observations indexed
$n = 1, \dots, N$ (each an $(a_n, v_n)$ pair with $v_n - a_n + 1 = L$). Write $P_n$ for that
observation's last observed paid amount, $g^{(c)}_n$ and $g^{(I)}_n$ for its (point-in-time)
company and industry age-to-lag factors, and $A_n$ for its true eventual (lag-10) paid amount
from the square. The blended estimate and its error as a function of the blend weight $r$ are:
\begin{align}
\widehat{P}_n(r) &= P_n \bigl[ r\, g^{(c)}_n + (1-r)\, g^{(I)}_n \bigr] \notag \\
e_n(r) &\equiv \widehat{P}_n(r) - A_n
       = \underbrace{\bigl(P_n g^{(I)}_n - A_n\bigr)}_{b_n} \;+\; r\underbrace{\Bigl(P_n\bigl(g^{(c)}_n - g^{(I)}_n\bigr)\Bigr)}_{x_n}
       = b_n + r\, x_n,
\label{eq:affine-error}
\end{align}
where $b_n$ is the error at $r=0$ (pure industry) and $x_n$ is the sensitivity of the error to
$r$ -- the gap between the company-only and industry-only point estimates. Every error is
therefore an affine (linear) function of the single scalar $r$, which is what makes both loss
functions below exactly, rather than only numerically, optimizable.

\subsection{RMSE: Closed-Form Quadratic Minimizer}

Since $e_n(r) = b_n + r x_n$ is affine, $\sum_n e_n(r)^2$ is a quadratic, upward-opening
(convex) parabola in $r$:
\begin{equation}
S(r) = \sum_{n=1}^{N} (b_n + r x_n)^2 = \sum_n b_n^2 \;+\; 2r\sum_n x_n b_n \;+\; r^2 \sum_n x_n^2.
\end{equation}
Differentiating and setting to zero:
\begin{equation}
\frac{dS}{dr} = 2\sum_n x_n b_n + 2r\sum_n x_n^2 = 0
\;\;\Longrightarrow\;\;
r^\star = -\frac{\sum_n x_n b_n}{\sum_n x_n^2},
\label{eq:rmse-star}
\end{equation}
with $\dfrac{d^2S}{dr^2} = 2\sum_n x_n^2 \ge 0$ confirming $r^\star$ is a global minimum whenever
$\sum_n x_n^2 > 0$ (if $\sum_n x_n^2 = 0$, every observation has $g^{(c)}_n = g^{(I)}_n$ and the
loss does not depend on $r$ at all). Since $r$ must remain a valid convex-combination weight, the
final RMSE-optimal ratio clips $r^\star$ to the unit interval:
\begin{equation}
\rrmse_{c,L} = \operatorname{clip}(r^\star,\,0,\,1), \qquad
\text{RMSE}\bigl(\rrmse_{c,L}\bigr) = \sqrt{\tfrac{1}{N}\textstyle\sum_n \bigl(b_n + \rrmse_{c,L}\, x_n\bigr)^2}.
\label{eq:rmse-clip}
\end{equation}

\subsection{MAE: Convex Piecewise-Linear Minimizer}

$\text{MAE}(r) = \tfrac{1}{N}\sum_n |b_n + r x_n|$ is a sum of $N$ convex, piecewise-linear
(``V-shaped'') functions of $r$ -- each term $|b_n + r x_n|$ has its single non-differentiable
kink (and, taken alone, its minimum) at $r = -b_n/x_n$ (undefined, i.e. constant in $r$, when
$x_n=0$). A finite sum of convex piecewise-linear functions is itself convex and piecewise-linear,
and the slope of a convex piecewise-linear function is non-decreasing in $r$; consequently its
minimum over an interval is attained either at one of the constituent kinks or at a boundary of
the interval. This yields an exact, finite candidate set:
\begin{equation}
\mathcal{C} = \{0,\,1\} \;\cup\; \bigl\{\operatorname{clip}(-b_n/x_n,\,0,\,1) : x_n \ne 0\bigr\},
\label{eq:mae-candidates}
\end{equation}
and the exact minimizer is found by direct evaluation:
\begin{equation}
\rmae_{c,L} = \argmin_{r\, \in\, \mathcal{C}} \ \frac{1}{N}\sum_{n=1}^N \bigl|b_n + r x_n\bigr|.
\label{eq:mae-argmin}
\end{equation}
No grid search or iterative solver is used for either loss function; both are exact given
Eqs.~\eqref{eq:rmse-clip} and \eqref{eq:mae-candidates}--\eqref{eq:mae-argmin}.

\subsection{Fully Worked Numerical Example}
\label{sec:worked-example}

Company 43, development lag $L=9$, has $N=2$ pooled observations:

\begin{table}[h]
\centering
\caption{Company 43, lag 9: the two pooled walk-forward observations}
\begin{tabular}{lrrrrrr}
\toprule
$a_n$ & $v_n$ & $P_n$ & $g^{(c)}_n$ & $g^{(I)}_n$ & $A_n$ & \\
\midrule
1998 & 2006 & 39{,}876 & 1.000000 & 1.000000 & 39{,}896 & \\
1999 & 2007 & 45{,}090 & 1.000502 & 1.001645 & 45{,}092 & \\
\bottomrule
\end{tabular}
\end{table}

Applying Eq.~\eqref{eq:affine-error}:
\begin{align*}
x_1 &= 39{,}876\,(1.000000 - 1.000000) = 0, &
b_1 &= 39{,}876(1.000000) - 39{,}896 = -20.000,\\
x_2 &= 45{,}090\,(1.000502 - 1.001645) = -51.549, &
b_2 &= 45{,}090(1.001645) - 45{,}092 = 72.164.
\end{align*}
\textbf{RMSE:} $\sum x_n b_n = (0)(-20.000)+(-51.549)(72.164) = -3720.02$;
$\sum x_n^2 = 0^2 + (-51.549)^2 = 2657.33$; so
$r^\star = -(-3720.02)/2657.33 = 1.3999$, clipped to $\rrmse_{43,9}=1.0$, giving
$\text{RMSE}(1.0) = \sqrt{\tfrac{1}{2}\bigl[(-20.0)^2 + (72.164-51.549)^2\bigr]} = \sqrt{412.51} = 20.310$.

\textbf{MAE:} the candidate set is $\{0,1\}$ (observation 1 contributes no kink since $x_1=0$;
observation 2's kink at $-b_2/x_2 = 1.400$ clips to the existing candidate 1). Evaluating:
$\text{MAE}(0) = \tfrac12(|{-20.0}|+|72.164|) = 46.082$;
$\text{MAE}(1) = \tfrac12(|{-20.0}|+|72.164-51.549|) = \tfrac12(20.0+20.615)=20.308$.
So $\rmae_{43,9} = 1.0$ as well, with $\text{MAE}=20.308$. In this instance both loss functions
agree that the company's own pattern should be fully trusted at this lag -- driven almost entirely
by observation 2, since observation 1 is uninformative about $r$ ($x_1=0$: company and industry
factors coincided exactly that year).

%==============================================================================
\section{Single-Company Case Studies}
\label{sec:casestudies}
%==============================================================================

\subsection{Company 43: Lag-to-Lag Instability of the Fitted Ratio}

Company 43's fitted RMSE-optimal ratio by lag illustrates that the fit is rarely a smooth function
of $L$:

\begin{table}[h]
\centering
\caption{Company 43: RMSE-optimal ratio and backing sample size, by development lag}
\begin{tabular}{lrrrrrrrrrr}
\toprule
Lag $L$ & 1 & 2 & 3 & 4 & 5 & 6 & 7 & 8 & 9 & 10 \\
\midrule
$N$ & 9 & 9 & 8 & 7 & 6 & 5 & 4 & 3 & 2 & 1 \\
$\rrmse$ & 0.664 & 0.510 & 0.153 & 0.488 & 1.000 & 0.902 & 0.504 & 1.000 & 1.000 & --- \\
\bottomrule
\end{tabular}
\end{table}

The ratio bounces between the interior and both boundaries as $L$ decreases -- symptomatic of the
declining sample size $N$ toward the higher lags (and toward $L=1$'s inherently thinner support):
with as few as 2--3 observations, $r^\star$ in Eq.~\eqref{eq:rmse-star} is determined almost
entirely by whichever single historical year had the largest residual, and RMSE loss (quadratic)
amplifies that single year's influence more than MAE loss (linear) would.

\subsection{Company 40568: The Convex-Combination Bound}
\label{sec:convex-bound}

A structural limit on how close \emph{any} blend can get to the actual outcome follows directly
from Eq.~\eqref{eq:blend-factor}: because $\widehat{P}^{\text{blend}} = r\,\widehat{P}^{(c)} +
(1-r)\,\widehat{P}^{(I)}$ is a convex combination (weights in $[0,1]$ summing to 1) of the two
single-source estimates, it must satisfy
\begin{equation}
\min\bigl(\widehat{P}^{(c)}, \widehat{P}^{(I)}\bigr) \;\le\; \widehat{P}^{\text{blend}} \;\le\;
\max\bigl(\widehat{P}^{(c)}, \widehat{P}^{(I)}\bigr) \qquad \text{for every } r \in [0,1].
\label{eq:convex-bound}
\end{equation}
Whenever the actual outcome $A$ falls outside this interval, \emph{no} choice of $r$ can bring
the blended estimate closer than the nearer of the two boundary endpoints -- this is not a
deficiency of the optimizer but a geometric fact about weighted averages. For company 40568,
accident year 2002 ($L=6$): the company-only estimate is $1374.73$, the industry-only estimate is
$1399.07$, and the true actual paid is $1472.00$ -- outside $[1374.73,1399.07]$ on the high side.
The RMSE-optimal fit correctly identifies $r=0$ (pure industry, distance $73.0$) as strictly
better than $r=1$ (pure company, distance $97.3$), which is the closest \emph{any} value of $r$
can get. The same pattern recurs at accident year 2005 ($L=3$): company-only $=1373.50$,
industry-only $=1369.89$, actual $=1203$ -- again outside the interval, this time on the low side,
with $r=0$ again marginally closer. Both cases reflect a genuine single-year anomaly in the
company's realized claims experience that neither a company-specific nor an industry-wide smoothed
development pattern can anticipate; no credibility weighting scheme built from Eq.~\eqref{eq:blend-factor}
can correct for it.

\subsection{Company 41700 and the Thin-History Flag}
\label{sec:thin}

Company 41700 (Southern Group Ind Inc) has recorded private-passenger-auto data for accident
years 1998 and 1999 only -- a full $10\times2$ block of development lags is present, but only two
accident-year cohorts exist at all. Every development lag's backtest fit therefore draws on at
most $N=2$ observations (and $N=1$ at the two most extreme lags), with several lags' ratios
landing exactly at $0$ or $1$ purely because two points admit no interior optimum under a
convex/piecewise-linear objective unless the interior point exactly minimizes both residuals
jointly. As described in Section~\ref{sec:population}, this company is flagged
(\texttt{thin\_history\_flag=True}, $\texttt{num\_accident\_years}=2$) both in the saved results
data and, prominently, on the rendered comparison chart.

%==============================================================================
\section{Cross-Sectional Analysis: Fitted Credibility versus Company Size}
\label{sec:crosssection}
%==============================================================================

\subsection{Company Share of Industry}

For each pooled $(\text{company } c, \text{lag } L)$ row, a size proxy is computed as the mean,
across that row's $N$ backing observations, of the company's share of total industry paid losses
at the matching cross-section:
\begin{equation}
\text{share}_{c,L} = \frac{1}{N}\sum_{n=1}^{N} \frac{P_n^{(c)}}{P_n^{(I,\text{total})}},
\label{eq:share}
\end{equation}
where $P_n^{(I,\text{total})}$ is the \emph{industry-wide} (sum across all 143 companies) last
observed cumulative paid loss at the same accident year and lag as observation $n$. This is a
paid-loss-based size proxy, distinct from (though correlated with) the premium-based concentration
figures of Section~\ref{sec:population}.

\subsection{Pooling and Reliability Filter}

Every company with a saved weighting-by-lag file (112 companies, up to 10 rows each) contributes
to a single pooled table of $1{,}120$ $(\text{company},L)$ rows (Section~\ref{sec:walkforward}).
Rows with $\texttt{num\_observations} < 2$ are excluded as unreliable (an exact fit to one point
carries no information about generalization), leaving $n=1{,}006$ rows for correlation analysis.
Over this reliable subset, $\text{share}_{c,L}$ ranges from $4.7\times10^{-7}$ to $0.742$ (median
$3.79\times10^{-4}$), a span of nearly six orders of magnitude.

\subsection{Corner-Clustering of the Fitted Ratio}

Among the 1{,}006 reliable rows, $\rmae$ sits at exactly $0$ for $231$ rows and exactly $1$ for
$344$ rows ($575/1006 = 57.2\%$ at a boundary); $\rrmse$ sits at exactly $0$ for $233$ and exactly
$1$ for $293$ ($526/1006=52.3\%$). This is a direct consequence of the optimizers derived in
Section~\ref{sec:optimization}: with $N$ typically in the single digits, the unconstrained RMSE
optimum in Eq.~\eqref{eq:rmse-star} frequently falls outside $[0,1]$ and is clipped, and the MAE
candidate set of Eq.~\eqref{eq:mae-candidates} frequently has its minimum at a boundary rather
than an interior kink.

\subsection{Correlation Measures}

Three correlation statistics are reported, each making a different assumption about the
relationship being tested:

\begin{itemize}
\item \textbf{Pearson $r$} assumes an approximately linear relationship between the two variables
  and is highly sensitive to both extreme skew in the independent variable (here, share spans six
  orders of magnitude) and to a dependent variable dominated by two repeated boundary values
  (here, the corner-clustering above), since it is fundamentally a measure of \emph{linear}
  co-movement.
\item \textbf{Spearman $\rho$} is the Pearson correlation computed on ranks (with tie-averaging),
  and detects any \emph{monotonic} (not necessarily linear) relationship, making it robust to the
  extreme skew in share but only partially robust to heavy tying in $r$.
\item \textbf{Kendall's $\tau_b$} is a concordant/discordant-pair statistic explicitly corrected
  for ties (the ``$b$'' variant, \texttt{scipy}'s default), making it the most appropriate of the
  three given that over half of $r$'s values are exact ties at $0$ or $1$.
\end{itemize}

\subsection{Results}

\begin{table}[h]
\centering
\caption{Correlation between company share of industry paid losses and fitted credibility ratio ($n=1{,}006$)}
\begin{tabular}{lrrrrrr}
\toprule
Ratio & Pearson $r$ & $p$ & Spearman $\rho$ & $p$ & Kendall $\tau_b$ & $p$ \\
\midrule
$\rmae$  & 0.0239  & 0.449 & 0.1288 & $4.17\times10^{-5}$ & 0.0929 & $3.64\times10^{-5}$ \\
$\rrmse$ & $-$0.0015 & 0.963 & 0.1060 & $7.57\times10^{-4}$ & 0.0769 & $5.49\times10^{-4}$ \\
\bottomrule
\end{tabular}
\end{table}

The Pearson coefficients are statistically indistinguishable from zero for both loss functions.
The rank-based statistics, in contrast, are small in magnitude but highly statistically
significant: a genuine, if weak, monotonic relationship between company size and fitted
credibility exists and is detectable only by measures robust to corner-clustering and skew --
Pearson correlation alone would lead an analyst to conclude, incorrectly, that no relationship
exists at all.

\subsection{Decile-Binned Credibility Curve}

As a complementary, assumption-light view, companies are grouped into ten equal-frequency bins
(deciles) of $\text{share}_{c,L}$ via quantile binning, and the mean and median fitted ratio are
computed within each bin ($n \approx 100$--$101$ per bin):

\begin{table}[h]
\centering
\caption{Mean and median fitted ratio by company-size decile (1 = smallest 10\% by share, 10 = largest)}
\begin{tabular}{lrrrrrr}
\toprule
& \multicolumn{3}{c}{$\rmae$} & \multicolumn{3}{c}{$\rrmse$} \\
\cmidrule(lr){2-4} \cmidrule(lr){5-7}
Decile & mean & median & $n$ & mean & median & $n$ \\
\midrule
1  & 0.675 & 1.000 & 101 & 0.629 & 0.842 & 101 \\
2  & 0.514 & 0.555 & 101 & 0.532 & 0.657 & 101 \\
3  & 0.598 & 0.853 & 100 & 0.545 & 0.699 & 100 \\
4  & 0.532 & 0.655 & 101 & 0.516 & 0.554 & 101 \\
5  & 0.540 & 0.632 & 100 & 0.549 & 0.628 & 100 \\
6  & 0.586 & 0.662 & 101 & 0.554 & 0.654 & 101 \\
7  & 0.581 & 0.655 & 100 & 0.641 & 0.784 & 100 \\
8  & 0.663 & 0.914 & 101 & 0.664 & 0.894 & 101 \\
9  & 0.670 & 0.911 & 100 & 0.622 & 0.861 & 100 \\
10 & 0.816 & 1.000 & 101 & 0.776 & 0.949 & 101 \\
\bottomrule
\end{tabular}
\end{table}

The pattern is not a clean monotone increase: both metrics are elevated at the extremes (decile 1
and deciles 8--10) relative to the middle deciles (2--7), an approximately ``smile''-shaped curve,
though with local bumpiness (e.g., decile 3 partially rebounds within the low-share region).
More informative than the mean/median levels themselves is the \emph{gap} between them within a
decile: a decile whose median sits at an exact boundary ($0$ or $1$) while its mean sits well
below/above indicates a \emph{polarized} population -- most companies in that decile are fully
trusted or fully distrusted, with the mean simply averaging across the two extremes -- as opposed
to a decile where mean and median are close together (a genuinely middling, unpolarized group).
Under $\rmae$, the largest mean--median gaps occur at decile 1 ($|0.675-1.000|=0.325$), decile 3
($0.256$), decile 8 ($0.251$), and decile 9 ($0.240$); the smallest occur at decile 2 ($0.041$)
and decile 7 ($0.074$). Decile 1 (the smallest companies) and decile 10 (the largest) are both
elevated and both polarized, but the middle deciles are not uniformly the least polarized either --
this nuance is invisible to a single correlation coefficient and is precisely the reason this
decile view is reported alongside Table~4.

%==============================================================================
\section{Volatility Analysis}
\label{sec:volatility}
%==============================================================================

\subsection{Motivation}

Every model in this project (company-only, industry-only, or blended) is a smoothed multiplicative
extrapolation of historical development. Such an extrapolation can only track a claims process
that itself behaves in a reasonably stable, patterned way; it cannot anticipate idiosyncratic,
noise-driven swings. This motivates a direct measure of how ``volatile'' or ``zigzag'' each
company's realized loss experience is, independent of company size, to test whether volatility --
rather than size per se -- is the more fundamental driver of both prediction error and fitted
credibility, consistent with the process-variance basis of classical credibility theory (a higher
process variance requires more observations, or equivalently receives lower credibility, before
full credibility is warranted; see Appendix~\ref{app:buhlmann}).

\subsection{Definition}

For company $c$, let $A_{c,1}, \dots, A_{c,T}$ ($T\le10$) be its ultimate (lag-10, fully matured)
actual paid loss, one value per accident year, sorted by accident year -- the same figures used as
grading truth throughout this report. Define the year-over-year relative change
\begin{equation}
\gamma_{c,t} = \frac{A_{c,t} - A_{c,t-1}}{A_{c,t-1}}, \qquad t=2,\dots,T,
\label{eq:growth}
\end{equation}
and the volatility statistic as the sample standard deviation (Bessel-corrected, divisor $T-2$
degrees of freedom lost: one for differencing, one for the mean) of the $T-1$ growth values:
\begin{equation}
\text{Vol}_c = \sqrt{\frac{1}{T-2}\sum_{t=2}^{T} \bigl(\gamma_{c,t} - \bar\gamma_c\bigr)^2}, \qquad
\bar\gamma_c = \frac{1}{T-1}\sum_{t=2}^{T}\gamma_{c,t}.
\label{eq:volatility}
\end{equation}
Relative (percentage) change rather than absolute dollar change is used deliberately: company size
in this dataset spans nearly six orders of magnitude (Section~\ref{sec:population}), so an
absolute-dollar measure would trivially rank the largest companies as ``most volatile'' regardless
of their actual relative stability. Companies with fewer than 5 accident years of history, or
fewer than 3 finite growth observations (after excluding divide-by-zero cases), are excluded,
leaving $n=105$ companies.

\subsubsection{Worked Example: Company 40568}

Ultimate actual paid by accident year: $1296, 1362, 1224, 1757, 1472, 1447, 1581, 1203, 1336,
1461$ (1998--2007). Year-over-year growth: $+5.09\%, -10.13\%, +43.55\%, -16.22\%, -1.70\%,
+9.26\%, -23.91\%, +11.06\%, +9.36\%$. Mean growth $\bar\gamma = 2.93\%$; sample standard
deviation $\text{Vol} = 19.62\%$ -- i.e., growth swings by roughly $\pm20$ percentage points
around a near-zero average trend, the numerical signature of a ``zigzag'' rather than a smoothly
growing or shrinking book of business.

\subsection{Volatility versus Prediction Error}

For each of the 105 companies, mean absolute value of the company-only backtest's
\texttt{error\_percentage} (defined as $100\times(\widehat{P}_{a,10}-A_{a})/A_{a}$, averaged
across that company's accident years) is correlated against $\text{Vol}_c$:

\begin{table}[h]
\centering
\caption{Volatility vs.\ mean absolute company-only error percentage ($n=105$)}
\begin{tabular}{lrr}
\toprule
& coefficient & $p$-value \\
\midrule
Pearson $r$  & 0.328 & $6.36\times10^{-4}$ \\
Spearman $\rho$ & 0.479 & $2.38\times10^{-7}$ \\
\bottomrule
\end{tabular}
\end{table}

Both are strongly significant and positive: more volatile companies are, unsurprisingly, harder
to predict.

\subsection{Volatility versus Company Size}

Merging $\text{Vol}_c$ with each company's mean $\text{share}_{c,L}$ (averaged across its lags,
share $>0$, $n=105$) gives
\begin{equation}
\text{Spearman}\bigl(\text{share}_c,\ \text{Vol}_c\bigr) = -0.683, \qquad p = 9.56\times10^{-16},
\end{equation}
a strong negative relationship: larger companies have systematically smoother loss experience.
The two least volatile companies in the dataset are State Farm Mut Grp ($\text{Vol}=0.049$) and
United Services Automobile Assn Grp ($\text{Vol}=0.029$) -- the two largest companies by premium
(Section~\ref{sec:population}) -- consistent with a law-of-large-numbers explanation: a larger,
more diversified book of policies averages out individual-claim-level randomness, while a small
company's aggregate paid losses are dominated by whichever handful of claims happen to occur in a
given year.

\begin{table}[h]
\centering
\caption{Most and least volatile companies ($n=105$, $\ge5$ accident years)}
\begin{tabular}{lrrr}
\toprule
Company code & Accident years & $\text{Vol}_c$ & Mean $|$error\%$|$ \\
\midrule
\multicolumn{4}{l}{\textit{Most volatile}} \\
26808 & 10 & 7.919 & 17.81\% \\
32743 & 7  & 1.711 & 2.91\% \\
32387 & 10 & 1.577 & 4.00\% \\
10308 & 10 & 1.556 & 12.09\% \\
13528 & 10 & 1.003 & 5.13\% \\
\midrule
\multicolumn{4}{l}{\textit{Least volatile}} \\
2003 (USAA) & 10 & 0.029 & 1.41\% \\
1090 (KY Farm Bureau) & 10 & 0.046 & 0.48\% \\
1767 (State Farm) & 10 & 0.049 & 0.29\% \\
620 & 10 & 0.059 & 2.20\% \\
6947 (Tenn Farmers Mut) & 10 & 0.064 & 0.75\% \\
\bottomrule
\end{tabular}
\end{table}

\subsection{Volatility versus the Fitted Ratio Directly}

Merging $\text{Vol}_c$ (one value per company) onto the pooled $(\text{company},L)$ table
(restricted to the 105 companies with a defined volatility statistic, $\texttt{num\_observations}
\ge 2$, $n=945$ rows) and correlating directly against the fitted ratio itself:

\begin{table}[h]
\centering
\caption{Volatility vs.\ fitted ratio ($n=945$ rows, 105 companies)}
\begin{tabular}{lrrrrrr}
\toprule
Ratio & Pearson $r$ & $p$ & Spearman $\rho$ & $p$ & Kendall $\tau_b$ & $p$ \\
\midrule
$\rmae$  & $-$0.0749 & 0.0213 & $-$0.0447 & 0.170 & $-$0.0313 & 0.178 \\
$\rrmse$ & $-$0.0849 & 0.00905 & $-$0.0229 & 0.482 & $-$0.0161 & 0.483 \\
\bottomrule
\end{tabular}
\end{table}

The sign is in the direction credibility theory predicts (higher volatility, lower fitted
credibility) but the relationship is weak, and only the Pearson statistic clears conventional
significance thresholds -- the opposite significance pattern from the size analysis of
Section~\ref{sec:crosssection}, where the rank-based statistics were the significant ones. The
most likely explanation is a mismatch of granularity: $\text{Vol}_c$ is a single scalar per
company, whereas the fitted ratio varies substantially \emph{within} a company across development
lags (Section~\ref{sec:casestudies}, Table~3, shows company 43's $\rrmse$ ranging from $0.153$ to
$1.000$ across lags). Pairing one company-level number against many lag-level outcomes for that
same company leaves most of the ratio's within-company variance unexplained by a between-company
statistic, diluting the correlation relative to volatility's much stronger relationship with
\emph{error magnitude} (Table~7), which is measured at the same company-level granularity as
$\text{Vol}_c$ itself.

%==============================================================================
\section{Discussion: How the Data Source Shapes These Results}
\label{sec:discussion}
%==============================================================================

Three properties of the underlying data recur across every analysis above and plausibly explain
why the size-credibility relationship is real but weak rather than strong:

\begin{enumerate}
\item \textbf{Extreme concentration.} With 92\% of industry premium held by 10 companies
  (Section~\ref{sec:population}) and company share of paid losses spanning nearly six orders of
  magnitude (Section~\ref{sec:crosssection}), any linear-correlation measure (Pearson) is
  structurally disadvantaged; only rank-based or binned views recover the signal.
\item \textbf{Thin per-lag samples.} A third of all pooled (company, lag) rows are backed by three
  or fewer walk-forward observations (Table~2); with so few points, both exact optimizers
  (Section~\ref{sec:optimization}) are prone to landing at a boundary ($r=0$ or $r=1$) driven by a
  single influential historical year rather than a stable underlying pattern -- the corner-clustering
  documented in Section~\ref{sec:crosssection} is a small-sample artifact layered on top of
  whatever genuine credibility signal exists.
\item \textbf{Single line, single geography, single ten-year window.} Private passenger auto is a
  comparatively short-tailed, well-behaved line; the industry-wide pattern is already a reasonably
  good predictor for most companies, which mechanically compresses how much room company size has
  left to explain in the residual credibility decision.
\end{enumerate}

%==============================================================================
\section{Limitations and Suggested Extensions}
\label{sec:limitations}
%==============================================================================

\begin{itemize}
\item \textbf{Unregularized point estimation.} $\rmae_{c,L}$ and $\rrmse_{c,L}$ are exact optima
  of Eqs.~\eqref{eq:rmse-clip} and \eqref{eq:mae-argmin} with no penalty for small $N$; a natural
  extension is to shrink the fitted ratio toward a prior (e.g., $0.5$, or a value smoothed across
  neighboring lags) in proportion to $1/N$ or a similar precision-weighted rule, in the spirit of
  limited-fluctuation or Bühlmann-Straub credibility (Appendix~\ref{app:buhlmann}), rather than
  trusting the unconstrained closed-form optimum at face value when $N$ is small.
\item \textbf{Out-of-sample generalization gap.} The fitted ratio is estimated entirely from
  valuation years prior to (or including, at the oldest backing accident year of each lag) the
  target year and then applied unchanged to 2007; no refitting occurs at the target year itself.
  This is methodologically correct for backtesting purposes but means the target-year blended
  estimate can miss by construction whenever 2007's realized experience differs from the historical
  average the ratio was fit to -- illustrated numerically in Section~\ref{sec:convex-bound}.
\item \textbf{No formal parametric credibility comparison.} This project fits $r$ empirically
  rather than via a parametric variance-components model (Bühlmann or Bühlmann-Straub); a natural
  extension is to compute the classical credibility factor $Z$ (Appendix~\ref{app:buhlmann}) from
  the same walk-forward panel and compare it directly against $\rmae$/$\rrmse$.
\item \textbf{Single-country, single-line scope.} All conclusions are specific to U.S.\ private
  passenger auto, 1998--2007; generalization to longer-tailed lines (general liability, workers'
  compensation) or other geographies is untested and, per Section~\ref{sec:discussion}, plausibly
  different given how much this line's short-tailed behavior compresses the residual credibility
  signal.
\end{itemize}

%==============================================================================
\section{Conclusion}
%==============================================================================

An empirically-fit credibility blend between company and industry paid chain-ladder development
factors, backtested via a rigorous walk-forward design with $1{,}006$ reliable (company,
development-lag) observations across 112 companies, shows a statistically significant but weak
positive rank correlation between company size and fitted credibility (Spearman $\rho \approx
0.11$--$0.13$), undetectable by Pearson correlation alone due to heavy boundary-clustering of the
fitted ratio. A directly measured volatility statistic explains substantially more of the
variation in prediction error (Spearman $\rho=0.479$) and is strongly negatively associated with
company size (Spearman $\rho=-0.683$), suggesting volatility -- consistent with the process-variance
foundation of classical credibility theory -- is the more fundamental quantity, with size acting
largely as its (imperfect) proxy in this dataset. Volatility does not, however, predict the fitted
ratio's \emph{direction} with comparable strength once measured at matching (company, lag)
granularity, a gap attributable to substantial within-company variation of the fitted ratio across
development lags that a single company-level volatility statistic cannot capture.

\appendix
%==============================================================================
\section{Classical Bühlmann Credibility Formula (Reference)}
\label{app:buhlmann}
%==============================================================================

For contrast with the empirically-fit ratio used throughout this report, the classical linear
(Bühlmann) credibility factor for a risk observed over $n$ periods is
\begin{equation}
Z = \frac{n}{n+K}, \qquad K = \frac{\text{EPV}}{\text{VHM}},
\end{equation}
where $\text{EPV}$ (expected value of the process variance) measures the average within-risk
variance across the population, and $\text{VHM}$ (variance of the hypothetical means) measures how
much the true mean loss level varies \emph{between} risks. $Z\to1$ as $n\to\infty$ or as
$\text{VHM}/\text{EPV}\to\infty$ (risks are very heterogeneous, so a given risk's own data is
highly informative about its own true mean); $Z\to0$ as $\text{EPV}/\text{VHM}\to\infty$ (risks
are essentially homogeneous, so an individual risk's own data is mostly noise around a common
mean). This formula was not fit in this project -- the ratio $r_{c,L}$ is instead obtained by
direct out-of-sample error minimization (Section~\ref{sec:optimization}) -- but the qualitative
prediction is the same one tested empirically here: a company whose own experience is more
informative relative to population-level noise (lower process variance, i.e., lower $\text{Vol}_c$)
should receive a higher credibility weight, exactly the (weak, but correctly-signed) relationship
found in Section~\ref{sec:volatility}.

%==============================================================================
\section{Software and Reproducibility}
%==============================================================================

All computation was performed in Python (pandas, NumPy, SciPy) against a src-layout package,
\texttt{reserve\_lab}. Key modules: \texttt{manual\_chain\_ladder.py} (Eqs.
\ref{eq:selected-factor}--\ref{eq:reserve}); \texttt{analysis.py}
(\texttt{run\_walk\_forward\_backtests}, Section~\ref{sec:walkforward};
\texttt{\_fit\_ratio}, Section~\ref{sec:optimization}; \texttt{calculate\_best\_ratio},
Eq.~\ref{eq:share}; \texttt{analyze\_credibility\_vs\_size}, Section~\ref{sec:crosssection});
\texttt{build\_loss\_structures.py} (triangle/square construction, Section~\ref{sec:square});
\texttt{compare.py} (all visualization, including the thin-history warning of
Section~\ref{sec:thin}). Correlation statistics use \texttt{scipy.stats.pearsonr},
\texttt{spearmanr}, and \texttt{kendalltau} (tau-b variant). All figures in this report were
recomputed directly from the project's saved CSV outputs
(\texttt{data/processed/results/*.csv}) at the time of writing.

\end{document}
